%----------------------------------------------------------------------------------------
% Cover letter for: "Relational biological structure improves fine-mapping of causal
% GWAS variants under weak signal" — submission to Nature Genetics
%
% Format follows the conventions of the companion GraphMana / GraphPop cover letters
% from the same author group (April 2026), which themselves follow standard Springer
% Nature cover-letter style.
%----------------------------------------------------------------------------------------

\documentclass[11pt,a4paper]{article}

\usepackage[top=0.85in,bottom=0.85in,left=1in,right=1in]{geometry}
\usepackage{parskip}
\usepackage{microtype}
\usepackage{enumitem}
\usepackage[hidelinks,colorlinks=true,urlcolor=blue,linkcolor=blue]{hyperref}

\setlength{\parindent}{0pt}
\setlist[itemize]{leftmargin=1.6em,topsep=0.2em,itemsep=0.2em,parsep=0pt}

\begin{document}

%----------------------------------------------------------------------------------------
% Sender header
%----------------------------------------------------------------------------------------

\textbf{Jian-Feng Mao, Ph.D.}\\
Professor, Department of Plant Physiology\\
Ume\aa{} Plant Science Centre, Ume\aa{} University\\
SE-901 87 Ume\aa{}, Sweden\\
Email: \href{mailto:jianfeng.mao@umu.se}{jianfeng.mao@umu.se} \quad$\cdot$\quad ORCID: \href{https://orcid.org/0000-0000-0000-0000}{0000-0000-0000-0000}

\vspace{0.4em}
\noindent\rule{\textwidth}{0.4pt}

\begin{flushright}
28 April 2026
\end{flushright}

\vspace{0.6em}

The Editors\\
\textit{Nature Genetics}\\
Springer Nature\\
London, United Kingdom

\vspace{0.8em}

\textbf{Re: Submission of ``Relational biological structure improves fine-mapping of causal
GWAS variants under weak signal'' as an Article}

\vspace{0.4em}

Dear Editors,

We are submitting ``Relational biological structure improves fine-mapping of causal GWAS
variants under weak signal'' for consideration as an \textbf{Article} at \textit{Nature Genetics}.
Because the contribution is a methodological reframing of how multi-omics annotations should
enter fine-mapping --- backed by cross-ancestry biobank and cross-species empirical results that
include canonical-lipid-gene recovery in Pan-UK Biobank and crop-genomics-canonical-gene
recovery on the 3{,}000 Rice Genomes panel --- we believe \textit{Nature Genetics} is the
natural venue.

Genome-wide association studies have catalogued tens of thousands of trait-associated loci, but
linkage disequilibrium leaves the actual causal variant statistically indistinguishable from
correlated neighbours. Fine-mapping is the post-GWAS step that resolves this. Existing Bayesian
fine-mappers (SuSiE, FINEMAP, PolyFun, and the recent infinitesimal extensions reported by
Wu \emph{et al.}\ in this journal in 2026) have converged on a common architecture in which
multi-omics annotations enter as a per-variant scalar weight vector. That representation
decouples cleanly from the statistical machinery, but it also collapses the relational structure
that should be the most informative signal --- and constructing the weight vector itself depends
on a stratified-LD-score-regression pipeline that is feasible in human biobanks and effectively
unavailable elsewhere. Five concrete shortcomings follow.

\textbf{Challenges GraphGWAS addresses}
\begin{itemize}
\item \textbf{Flattened annotation priors} --- gene, tissue, pathway, regulatory-element and
protein--protein-interaction context that connects each variant to its biological neighbourhood
is collapsed to a single per-variant number; the relational structure that should break LD ties
is discarded before the inference runs.
\item \textbf{Annotation-pipeline bottleneck} --- the scalar prior requires stratified LD-score
regression over a curated functional-category catalogue (Baseline 2.2 in human uses 96
categories), which assumes a large GWAS, pre-curated categories, and an ancestry-matched LD
reference --- prerequisites met in human biobanks and absent in crop, livestock, and model-organism GWAS.
\item \textbf{Per-locus speed at biobank scale} --- runtimes of 1--3 seconds per locus for
vanilla SuSiE/FINEMAP, plus repeated reconstruction of each locus's multi-omics neighbourhood
from separate BED, GTF and TSV files, make annotation-aware biobank-scale fine-mapping
practically painful.
\item \textbf{Cross-ancestry brittleness} --- ancestry-private variants and substantially
different LD across ancestries are not handled by single-population fine-mappers; ancestry-by-
ancestry pipelines are typically rebuilt by downstream users.
\item \textbf{No cross-species portability} --- the Baseline-functional-category framework is
human-specific; a researcher fine-mapping yeast, \emph{Arabidopsis}, or rice GWAS must rebuild
the pipeline from scratch with annotations that are sparse, fragmented and not pre-harmonised.
\end{itemize}

We address this with a typed factor graph that carries the multi-omics neighbourhood through
inference rather than collapsing it to a scalar, and with two fine-mapping methods that exploit
that representation directly.

\textbf{Key novelties of this work}
\begin{itemize}
\item \textbf{Hierarchical belief propagation (HBP)} --- message passing on a three-layer
variant--gene--pathway factor graph with PPI coupling, with proved geometric convergence via a
Banach contraction argument (Theorem 2). The full multi-omics neighbourhood enters every
update.
\item \textbf{Graph-augmented fine-mapping (GAFM)} --- an annotation-adaptive Bayesian
complement that LD-deconvolves z-statistics and combines them with a multi-layer functional
score, with a proof that the causal variant ranks first under mild LD-decay assumptions
(Theorem 3). At weak signal with tissue-specific eQTL priors, GAFM beats SuSiE 27 to 2
head-to-head ($p < 10^{-5}$).
\item \textbf{Speed at parity accuracy} --- on a uniform 6-method $\times$ 3-scenario benchmark
on 1000 Genomes chromosome~22 (90 simulations), HBP and GAFM match the rank-\#1 rate of SuSiE,
FINEMAP, SuSiE-inf and FINEMAP-inf within three percentage points at 5--40$\times$ the speed
(median 0.07--0.12\,s per locus versus 0.22--3.1\,s for the baselines).
\item \textbf{Pan-UK Biobank sumstats-streaming entry path} --- per-locus tabix-over-HTTPS
fetching and on-the-fly GRCh37$\to$GRCh38 lifting, applied unchanged to BMI, LDL,
triglycerides and height across four Pan-UKB ancestries (EUR $N{=}420{,}531$;
CSA, AFR, EAS) at single-variant resolution.
\item \textbf{Cross-species portability without algorithmic change} --- the same code
fine-maps the 1011 Yeast Genomes panel (35 traits, 245 leads), the 1001 \emph{Arabidopsis}
Genomes panel (14 traits, 54 leads) and the 3{,}000 Rice Genomes panel (18 IRRI ordinal
traits, 72 leads + 4 continuous grain weight/shape traits, 41 fine-mapped leads). On rice,
24 of 72 IRRI-trait leads fall within 250\,kb of a Ren \emph{et al.}\ (2023) catalogued
causal gene; on the panel-matched grain weight + shape pass, \textbf{20 of 21 stable QTNs
reported by Niu \emph{et al.}\ 2021 \emph{BMC Genomics} are recovered at GWAS lead
resolution}, and at the variant-level metric the mixture-prior posterior reweightings
of GAFM and HBP (GAFM-MX, HBP-MX) and their ensemble (ENS) each reach 57\% recovery
with a top-1-PIP exact-position rate of \textbf{47.6\%
(10 of 21) --- the highest of any method tested, beating SuSiE (28.6\%) and SBayesRC (14.3\%)
on this metric while remaining 200--700$\times$ faster per locus than SuSiE}. SBayesRC
itself, run on a custom 3kRG eigen-decomposed LD reference we built from scratch
(GCTB 2.05beta + R LDstep1--4, 23 targeted $\pm$250\,kb blocks, ${\sim}$10\,min compute),
achieves 95\% Tier-1 recovery; the mixture-prior reweighting demonstrates that this
gap can be substantially closed by importing the same effect-size mixture prior into
GAFM/HBP as a posterior reweighting step.
\end{itemize}

A direct-intervention experiment establishes the central methodological claim that the
typed-graph representation, not annotation count, is what lets the graph break LD ties.
Replacing a uniform GTEx + STRING human cache with the same cache plus a single coarse
nine-class ENCODE regulatory-element flag --- no other change to GWAS sumstats, LD references,
or the fine-mapping algorithm --- raises the HBP-tighter-than-GAFM rate from 41\% to 92\%
across all 692 leads in the four-species replication. Density alone is therefore
insufficient even at $10^7$-variant scale; the operative quantity is per-variant
discriminative information, which the typed-graph representation exposes directly to inference.

Graph-native fine-mapping has practical implications beyond benchmarks: it removes the
annotation-curation bottleneck that has prevented systematic fine-mapping in non-model
species --- crops, livestock, and model organisms. The 3{,}000 Rice Genomes recovery of
\emph{Rc}, \emph{TGW6}, \emph{An-1}, \emph{BG1}, \emph{GAD1/RAE2} and \emph{OsER1} on the
IRRI ordinal-trait pass, together with the 20-of-21-stable-QTN recovery against the
panel-matched Niu \emph{et al.}\ 2021 grain weight/shape catalogue, all from the same code
base that produces single-variant credible sets at \emph{LDLR} and \emph{LPL} on Pan-UK
Biobank, is a concrete demonstration of this portability on the canonical crop-genomics
benchmark species.

A companion manuscript on graph-native higher-order epistasis (LD-pruned co-occurrence
pairwise-epistasis; LPCE) is in preparation; we will alert the journal upon its submission so
the two papers can be considered together if the editorial team finds it useful. Two further
methodological papers from the same author group are being submitted concurrently to
\textit{Nature Methods} on the underlying data-management (GraphMana) and population-genetics
computation (GraphPop) substrate; those papers do not overlap with the fine-mapping
contribution reported here.

\textbf{Code and data availability.}
Source code is released under the MIT licence at
\href{https://github.com/jfmao/GraphGWAS}{https://github.com/jfmao/GraphGWAS}. A staged Zenodo
deposit (graph-database dumps for human 1000 Genomes with multi-omics annotations and yeast
1011, multi-omics graph caches with the heterogeneity-intervention extension for all four
species, benchmark JSON outputs, Pan-UK Biobank per-locus fine-mapping outputs, and the full
3{,}000 Rice Genomes IRRI 18-trait summary statistics) is ready for upload and will receive a
DOI on acceptance.

\textbf{Suggested and non-preferred reviewers.}
Suggested and non-preferred reviewers (with affiliations, contact details, and brief rationale)
are provided in the submission portal ``Reviewer suggestions'' fields. We suggest reviewers
with combined expertise in statistical fine-mapping methodology, multi-omics annotation
integration, biobank-scale and cross-ancestry GWAS, and crop / model-organism genomics; we
have flagged close collaborators and recent co-authors as non-preferred.

No part of this work has been published or submitted elsewhere. All authors have read and
approved the manuscript. The authors declare no competing interests.

\vspace{0.6em}

Sincerely,

\vspace{1.4em}

\textbf{Jian-Feng Mao, Ph.D.}\\
Corresponding author, on behalf of all authors (Estaji, Zhao, Chen, Nie, and Mao)\\
Ume\aa{} Plant Science Centre, Department of Plant Physiology, Ume\aa{} University\\
SE-901 87 Ume\aa{}, Sweden\\
Email: \href{mailto:jianfeng.mao@umu.se}{jianfeng.mao@umu.se}

\vspace{1em}

\textit{Enclosures: manuscript, supplementary information, reviewer suggestions (submission portal).}

\end{document}
